\documentclass[a4paper, 12pt]{report}

% =======================
% IMPORTATION DES PACKAGES
% =======================

\usepackage[T1]{fontenc}
\usepackage[utf8]{inputenc}
\usepackage[margin=2cm]{geometry}
\usepackage{lmodern}
\usepackage{theorem}
\usepackage{amsfonts, amsmath, amssymb, dsfont}
\usepackage[french]{babel}

\renewcommand{\familydefault}{\sfdefault}

\pagestyle{headings}

\title{TOPOLOGIE}
\date{\today}
\author{}

\theoremstyle{break}
\theoremheaderfont{\scshape}
\theorembodyfont{\upshape}

\newtheorem{defin}{Définition}[section]
\newtheorem{prop}[defin]{Proposition}
\newtheorem{theo}[defin]{Théorème}
\newtheorem{exe}[defin]{Exemple}

\begin{document}

\maketitle

\textbf{PRÉCISION}\\
Je tient à préciser que ceci n' est pas un book mais juste un recueille des notes pour approfondir mes connaissances en mathématiques. 
Tout en développant mes compétences en rédaction scientifique.
Je le reformule avec mes propres mots, sauf les sous-titres.
Si l'envie vous tente, vous pouvez consulter mes notes.

\tableofcontents
% ==========
% CHAPITRE 1 
% ==========

\chapter{ESPACE TOPOLOGIQUE}

\section{Espace topologique et propriété}
\begin{defin}[Topologie]
Soit $X$ un espace vectoriel normé et $\tau$ un famille d'ouvert qui contient les partis de $X$.
On dit que $\tau$ est un topologie si elle vérifie les conditions suivants:
\begin{enumerate}
	\item $\emptyset \in \tau$  et $X \in \tau$
	\item Si $(O_{i})_{i \in \mathbb{N}} \in \tau$, on a $\bigcup_{i \in \mathbb{N}} O_{i} \in \tau$
	\item Si $(O_{i})_{i \in \mathbb{N}} \in \tau$, on a $\bigcap_{i \in \mathbb{N}} O_{i} \in \tau$
\end{enumerate}
\end{defin}
Le couple $(X, \tau)$ est un espace topologique. $\tau$ est un ouvert.

\begin{exe}
	\begin{enumerate}
		\item $\tau = P(X)$ est une topologie, elle est appelé topologie discrète.
		\item $\tau = \{0, X\}$ est une topologie appelé topologie grossière.
		\item $\tau = \{ O \subset X, X \backslash O \mbox{ finie}\} \cap \{ \emptyset \}$ est une topologie appelé topologie cofinie.
	\end{enumerate}
\end{exe}
\textit{Preuve.}
\begin{enumerate}
	\item Soit $(X, \tau)$ un espace topologique et $P(X)$ l'ensemble des parties de X.\\
	$\emptyset$ et X sont des parties de X, ainsi $\emptyset$ et $X \in \tau$.\\
	Soit $O_{1}, O_{2} \in \tau$, Si $O_{1} \cap O_{2} = \emptyset$ alors $O_{1} \cap O_{2} \in \tau$, sinon si $O_{1} \cap O_{2} \not= \emptyset$ et si $O_{1} \subset O_{2}$ alors $O_{1} \cap O_{2} = O_{1} \in \tau$ de même si $O_{2} \subset O_{1}$.\\
	Soit $(O_{i})_{i \in \mathbb{N}} \in \tau$, $\bigcup_{i \in \mathbb{N}} O_{i} \in \tau$.\\
	Ainsi $\tau = P(X)$ est bien une topologie sur $X$.

	\item Soit $(X, \tau)$ un espace topologique.\\
	Par définition de $\tau$, $\emptyset$ et X $\in \tau$.\\
	On a $\emptyset \cap X = \emptyset \in \tau$, ainsi elle est stable par intersection.\\
	On a $\emptyset \cup X = X \in \tau$, elle est stable par union.\\
	Donc $\tau = \{ \emptyset, X\}$ est bien une topologie sur $X$.
	\item Soit $(X, \tau)$ un espace topologique.\\
	Par définition de $\tau$, $\emptyset$ et $X \in \tau$ car $X\backslash \emptyset = X \in \tau$ et $X\backslash X = \emptyset \in \tau$.\\
	Soit $O_{1}, O_{2} \in \tau$, on sait que $X\backslash O_{1}$ et $X \backslash O_{2}$ sont de cardinal finie. Donc on a $O_{1} \cap O_{2}$ est aussi de cardinal finie. Ainsi $O_{1} \cap O_{2} \in \tau$.\\
		Soit $(O_{i})_{i \in \mathbb{N}} \in \tau$, on a $\bigcup_{i \in \mathbb{N}} O_{i} = X \backslash (\bigcap_{i \in \mathbb{N}} O_{i})$ est finie, donc $\bigcup_{i \in \mathbb{N}} O_{i} \in \tau$.\\
	On conclut alors que $\tau = \{ O \subset X, X \backslash O \mbox{ finie}\} \cap \{ \emptyset \}$ est une topologie sur X. \marginpar{$\square$}
\end{enumerate}

\begin{defin}[Affinités de la topologie]
      Soit $(X, \tau_{1})$, $(X, \tau_{2})$, des espaces topologiques. On dit que $\tau_{1}$ est moins fine que $\tau_{2}$ ou que $\tau_{2}$ plus fine que $\tau_{1}$ si $\tau_{1} \subseteq \tau_{2}$.
\end{defin}

\begin{defin}[Base de la topologie]
      Soit $(X, \tau)$ un espace topologique et $\mathcal{B} \subseteq \tau$. On dit que $\mathcal{B}$ est une base topologique si pour $O \in \tau$, il existe $(U_{i})_{i \in \mathbb{N}} \in \mathcal{B}\mbox{ tel que } O = \displaystyle \bigcup_{i \in \mathbb{N}} U_{i}$. 
\end{defin}

\begin{theo}
      Soit $(X, \tau)$ un espace topologique et $\mathcal{B} \subseteq \tau$. Lasse :
      \begin{enumerate}
               \item $\mathcal{B}$ est une base topologique de $\tau$,
               \item Pour chaque $U \in \tau$ et $x \in U$, on peut trouver $O \in \mathcal{B}$ tel que $x \in O \subseteq U$.
      \end{enumerate}
\end {theo}

\begin{defin}[Fermé]
      Soit $X$ un sous ensemble non vide. On dit que le sous ensemble $F$ de $X$ est fermé si son complémentaire est ouvert, c'est à dire $X\backslash F$ est ouvert alors $F$ est fermé.
\end{defin}

\begin{defin}[Ouvert]
      Soit $X$ un ensemble non vide. On dit que $U$ est un ouvert de $X$ si pour tout $x \in U, \; \mbox{ il existe }\epsilon >0$  tel que $B(x,\epsilon) \subseteq U$
\end{defin}

\begin{theo}
     \begin{enumerate}
           \item $\emptyset$ et $X$ sont à la fois ouvert et fermé,
           \item L'intersection de tous ensembles fermé (finie ou non) est fermé,
           \item L'union de tous ensembles fermé est un fermé,
           \item L'intersection de tous ensembles ouvert est un ouvert,
           \item L'union de tous ensembles ouvert (finie ou non) est un ouvert.
     \end{enumerate}
\end{theo}

\begin{defin}[Voisinage]
     Soit $(X,\tau)$ un espace topologique et $V$ un sous ensemble de X.
     \begin{itemize}
            \item On dit que $V$ est voisinage de x dans $X$, s'il existe $U \in \tau$ tel que $x \in U \subset V$. On note $\mathcal{V}(x)$, l'ensemble de tous les voisinages de x.
            \item De manière analogue, on dit que $V$ est voisinage de $A$ dans $X$, s'il existe $U \in \tau$ tel que $A \subset U \subset V$. On note $\mathcal{V}(A)$, l'ensemble de tous les voisinages de x.
     \end{itemize}
\end{defin} 

\begin{theo}
Soit $(X, \tau)$ un espace topologique.
\begin{enumerate}
      \item L'ensemble $\mathcal{V}(x)$ n'est jamais vide,
      \item Soit $(V_{i})_{i \in \mathbb{N}} \in \mathcal{V}(x)$, alors $\displaystyle\bigcap_{i \in \mathbb{N}} V_{i} \in \mathcal{V}(x)$,
      \item Soit $V \in \mathcal{V}(x)$, et $W \subset X$ alors $W \in \mathcal{V}(x)$.
\end{enumerate}
\end{theo}

\begin{prop}
 Soit $(X, \tau)$ un espace topologique et $V$ un sous ensemble de X. Lasse :
 \begin{enumerate}
         \item $V$ est un ouvert de $X$,
         \item $V$ est une voisinage de chacun de ses points c'est à dire $V \in \mathcal{V}(x)$.
 \end{enumerate}
\end{prop}

% ==========
% CHAPITRE 2
% ==========

\chapter{LIMITE ET CONTINUITÉ}

\section{Filtre et ultrafiltre}
\begin{defin}[Relation d'ordre]
       Soit X un ensemble non vide. On dit que la relation $\preccurlyeq$ est une relation d'ordre si elle vérifie les propriété suivants:
       \begin{itemize}
              \item Reflexivité: pour tout $x \in X, \; x \preccurlyeq x$
              \item Anti-symetrie:  pour tout $x, y \in X, x \preccurlyeq y, y \preccurlyeq x$, alors $x=y$,
              \item Transitivité : pour tout $x,y,z \in X, x \preccurlyeq y \mbox{ et } y \preccurlyeq z$ alors $x \preccurlyeq z$.
       \end{itemize}
       L'ensemble $(A, \preccurlyeq)$ est appelé ensemble ordonnée.
\end{defin}

\begin{defin}[Filtrant à droite]
        On dit que $(A, \preccurlyeq)$ est un ensemble dirigé ou un filtrant à droite si pour tout $\alpha, \beta \in (A \times A)$, il existe $\gamma \in A$ tel que $\alpha \preccurlyeq \gamma$ et $\beta \preccurlyeq \gamma$.
\end{defin}

\begin{defin}[Ensemble cofinal]
      Soit $(A, \preccurlyeq)$ un ensemble dirigé. On dit que $B \subseteq A$ est cofinal si pour tout $\alpha \in A$, on peut trouver $\beta \in B$ tel que $\alpha \preccurlyeq \beta$.
\end{defin}

\begin{defin}[Filtre]
       On dit que $\mathcal{F}$ est un filtre si elle vérifie les conditions suivants:
       \begin{itemize}
             \item $\emptyset \not\in \mathcal{F}$ et $X \in \mathcal{F}$,
             \item $A_{1}, A_{2}, \cdots, A_{n} \in \mathcal{F}$, alors $\displaystyle \bigcap_{i \in \mathbb{N}} A_{i} \in \mathcal{F}$,
             \item $A \in \mathcal{F}, \; B \in \mathcal{P}(X)$ tel que $A \subseteq B$, alors $B \in \mathcal{F}$.
       \end{itemize}
\end{defin}

\begin{defin}
Soit $\mathcal{F}$, un filtre
\begin{enumerate}
      \item Une suite généralisé est une suite $(x_{\alpha})_{\alpha \in A}$ où $A$ est un ensemble dirigé.
\end{enumerate}
\end{defin}
\end{document}
